% Generated draft: requires mathematical and source review.
\documentclass[11pt]{article}
\usepackage[margin=25mm]{geometry}
\usepackage{fontspec}
\setmainfont{Latin Modern Roman}
\usepackage{amsmath,amssymb,mathtools,unicode-math}
\setmathfont{Latin Modern Math}
\usepackage{enumitem,needspace,xurl,hyperref}
\setlength{\parindent}{0pt}
\setlength{\parskip}{5pt}
\setlength{\emergencystretch}{3em}
\newcommand{\N}{\mathbb{N}}
\newcommand{\Z}{\mathbb{Z}}
\newcommand{\Q}{\mathbb{Q}}
\newcommand{\R}{\mathbb{R}}
\newcommand{\C}{\mathbb{C}}
\newcommand{\F}{\mathbb{F}}
\newcommand{\A}{\mathbb{A}}
\newcommand{\PP}{\mathbb{P}}
\newcommand{\E}{\mathbb{E}}
\newcommand{\eps}{\varepsilon}
\newcommand{\dd}{\,\mathrm{d}}
\begin{document}
\section*{UnsolvedMath ID: 800139}
\textbf{Source problem code:} catalogue entry
\subsection{Definitions and mathematical statement}
The supplied record asks for bounds on Skewes's numbers of other highly biased prime races, conjecturally or provably. It mentions related investigations in number fields by Fiorilli and Jouve [FJ] and Bailleul [Bai].

For an intended race \(\mathcal R\), write \(\mathbf p(\mathcal R)\) for its defining parameters and \(S(\mathcal R)\) for its Skewes number, once these objects have been specified. The requested task is:
\begin{quote}
Obtain explicit bounds for \(S(\mathcal R)\) in terms of \(\mathbf p(\mathcal R)\), stating all hypotheses for each conjectural or proved bound.
\end{quote}
This is only a formulation schema. The supplied excerpt does not specify the races, their counting functions, their parameter ranges, the meaning of ``highly biased,'' or the convention defining a Skewes number. In particular, it does not determine the direction of the inequality whose first reversal is sought, whether equality counts as a reversal, the starting threshold, or the treatment of races with no reversal. These choices cannot be supplied silently, so the excerpt does not determine a single precise bounding problem.

The record also contains a separately attributed passage from Alexandre Bailleul, Universit\'e Paris--Saclay, concerning Chebyshev's bias over function fields. To explain the distinction in that passage, for a selected finite family of real zero arguments \(\theta_1,\ldots,\theta_m\), linear independence over \(\Q\) means
\[
\left(\sum_{j=1}^{m}q_j\theta_j=0,\quad q_1,\ldots,q_m\in\Q\right)
\Longrightarrow q_1=\cdots=q_m=0.
\]
Linear independence together with \(\pi\) means
\[
\left(q_0\pi+\sum_{j=1}^{m}q_j\theta_j=0,\quad q_0,\ldots,q_m\in\Q\right)
\Longrightarrow q_0=\cdots=q_m=0.
\]
For an infinite selected family, each condition is imposed on every finite subfamily. The source does not specify the relevant \(L\)-functions, which zero arguments are selected, their representatives, or how multiplicities are handled.

The supplied passage reports that, following Cha [Cha], function-field treatments typically assume independence together with \(\pi\), in addition to independence of the zero arguments themselves. It reports no known rational linear relations among the arguments themselves, while reporting known instances of an argument in \(\pi\Q\). It says these instances usually arise from special geometric conditions, such as supersingularity of elliptic curves or Jacobians that are not absolutely simple. These are supplied source qualifications, not independently verified assertions. No explicit question accompanies this passage, and its connection to the Skewes-number bounding request is unresolved.
\subsection{Short English statement}
The main request is to bound the first reversal thresholds called Skewes's numbers for highly biased prime races using the parameters defining those races. The excerpt leaves the races and threshold conventions unspecified. An additional function-field passage discusses independence of zero arguments, including independence together with \(\pi\), but supplies no explicit question.
\subsection{Sources}
\begin{itemize}
\item Local UnsolvedMath registry: supplied complete record with permanent ID \(800139\), source code \texttt{catalogue entry}, title ``Comparative Prime Number Theory --- Problem 8,'' in the AMR Open Problem Lists collection. The registry labels the record \texttt{open}; this is a supplied status, not an independently verified mathematical assessment.
\item Hamieh, Kadiri, Martin, and Ng, \emph{Comparative Prime Number Theory Problem List} (2024), Problem (\#8), \url{https://arxiv.org/abs/2407.03530}. This is the supplied primary-source attribution. The record reports access on 2026-10-16 and extraction by \texttt{arxiv\_latex}; neither access nor extraction has been independently checked here.
\item [FJ]: Fiorilli and Jouve, cited in the supplied excerpt for related highly biased races in number fields. No title, publication details, or separate URL is supplied.
\item [Bai]: Bailleul, cited in the supplied excerpt for related highly biased races in number fields. No title, publication details, or separate URL is supplied.
\item [Cha]: Cha, cited in the supplied excerpt concerning customary independence assumptions in function-field Chebyshev bias. No title, publication details, or separate URL is supplied.
\item The function-field passage is attributed in the supplied record to Alexandre Bailleul, Universit\'e Paris--Saclay. Its placement within the intended problem is unresolved.
\end{itemize}
\end{document}
